\section{Discussion, Limitations and Future Work}

This section discusses the design choices of DSphinx, their implications, and the remaining limitations of the proposed approach.

\smallskip
\noindent\textbf{Impact on route manipulation attacks.}
DSphinx aims to reduce a malicious client's ability to manipulate routing decisions. 
However, it does not completely eliminate the possibility of route manipulation attacks.
In conventional source-routing systems, clients are free to construct arbitrary routes and repeatedly force the inclusion or exclusion of specific nodes.
By contrast, DSphinx derives routes through a decentralized sampling procedure, preventing clients from directly controlling path selection. 
Although an adversary may repeatedly query the sampling mechanism until obtaining a desirable route, this is substantially more costly than arbitrary route construction.
Once a route is sampled, it may be reused multiple times by re-randomizing $\A{}$ and the per-hop shared secrets. 
Furthermore, repeated reuse of the same sampled route may create observable traffic patterns, making route-manipulation attacks easier to detect.

\smallskip
\noindent\textbf{Limitations and Future Work.}
DSphinx has few limitations. 
First, a deployment challenge in establishing a sufficiently large, diverse, and independent set of STTP operators. 
Second, it relies on a threshold trust assumption: compromise of at least $d$ STTPs breaks header confidentiality and enables full path reconstruction. 
Third, decentralizing route sampling introduces additional computation and communication overhead. 
Although modest (less than 50 ms in our evaluation) compared with the end-to-end latency of decentralized networks such as the Lightning Network and mixnets, these costs should be considered in latency-sensitive settings.

Finally, the current design focuses on \emph{random} route sampling, making it directly applicable to anonymous communication systems such as mixnets, where uniformly random paths are desirable. 
Extending the approach to routing algorithms based on optimization objectives, such as shortest-path, minimum-cost, or congestion-aware routing used in payment-channel networks like Lightning, remains an open research direction.
